% main.tex — The local structure of the random k-out graph
% Compile: tectonic main.tex  (all packages freely fetched by Tectonic)
\documentclass[11pt,a4paper]{article}
\usepackage[margin=2.6cm]{geometry}
\usepackage{amsmath,amssymb,amsthm}
\usepackage{booktabs}
\usepackage{xcolor}
\usepackage{url}
\usepackage[colorlinks=true,linkcolor=blue!50!black,citecolor=blue!50!black,urlcolor=blue!50!black]{hyperref}
\usepackage{microtype}

\newtheorem{theorem}{Theorem}[section]
\newtheorem{lemma}[theorem]{Lemma}
\newtheorem{corollary}[theorem]{Corollary}
\newtheorem{proposition}[theorem]{Proposition}
\newtheorem{conjecture}[theorem]{Conjecture}
\theoremstyle{definition}
\newtheorem*{remark}{Remark}
\theoremstyle{definition}
\newtheorem*{definition*}{Definition}

\title{Exact expected subgraph counts in the random $k$-out graph:\\
closed forms, a matched-density Erd\H{o}s--R\'enyi contrast,\\ and an orientation invariant}
\author{Rajveersinh Vishal Pardeshi\thanks{Independent researcher. No institutional
affiliation. Code, data, and the verification documents described in
Section~\ref{sec:repro} accompany this work in a public repository; see
Section~\ref{sec:repro} for the reproducibility statement.}}
\date{October 18, 2026}

\begin{document}
\maketitle

\begin{abstract}
In the \emph{random $k$-out graph} $G(n,k\text{-out})$ each of the $n$ vertices
independently selects a $k$-subset of the other $n-1$ vertices, and two vertices are
adjacent if either selected the other (symmetrized union). The model is a standard
$O(nk)$-edge alternative to the Erd\H{o}s--R\'enyi graph $G(n,p)$ and has been studied
chiefly through global observables (connectivity, perfect matchings, Hamilton cycles).
We develop its \emph{local structure}. Our first result is an exact closed formula,
valid for every fixed finite simple graph $H$ and all finite $(n,k)$, for the expected
number of copies of $H$ (Theorem~\ref{thm:general}); it relies on a per-chooser
avoidance probability that is exact hypergeometric and on chooser independence, and it
quantifies the \emph{within-chooser negative correlation}
$\Pr(v\in S_u,\, w\in S_u)=\frac{k(k-1)}{(n-1)(n-2)}$ that invalidates naive
$G(n,p)$-style shortcuts. We then compute: (i) the exact triangle closed form
$P(\triangle)=1-3A^{2}+3A^{2}B-B^{3}$ and its series, the limiting mean
$\mathbb E T \to (4k^{3}-3k^{2})/3$, the limiting variance
$\operatorname{Var}(T)\to(4k^{3}-3k^{2})/3$ (with exact $1/n$ terms), and the exact
$k=1$ specials $P(\triangle)=2/(n-1)^3$ (Theorems~\ref{thm:tri},
\ref{thm:variance}); (ii) a matched-density comparison with $G(n,p)$: the expected
triangle counts obey a ratio tending to $1-\frac{3}{4k}$, with the four-cycle ratio
tending to $1-\frac{1}{k}+\frac{1}{8k^{2}}$ (Proposition~\ref{prop:er}); (iii) an
exact cancellation theorem for stars $K_{1,m}$: the first non-zero term of the
fixed-vertices star probability is exactly of order $n^{-m}$, and
$\mathbb E[\,\#\text{ unlabelled }K_{1,m}\,] = \frac{c_m(k)}{m!}\,n\,(1+o(1))$ with
explicit $c_m(k)$ (Theorem~\ref{thm:stars}); and (iv) a general asymptotic
description: for any fixed $H$, $\mathbb E[\,\#\text{ labelled copies of
}H\,]\sim\Omega_k(H)\, n^{v-e}$, where
\[
\Omega_k(H)\;=\;\sum_{\sigma}\ \prod_{x\in V(H)}(k)_{t_x(\sigma)}
\]
is an explicit \emph{orientation invariant} summed over the assignments $\sigma$ of
each edge of $H$ to exactly one of its endpoints, $(k)_t$ being a falling factorial
(Theorem~\ref{thm:omega}). $\Omega_k(H)=0$ if and only if no assignment keeps every
vertex's owned degree at most $k$. All approximate statements above carry
machine-exact series constants. All closed forms are verified at exact precision:
symbolically (sympy), by exhaustive enumeration over \emph{all} choice profiles for
small $n$ (difference zero), and by Monte Carlo at large $n$ (all $|z|\le 1.43$). The
Poisson-type limit of the triangle count is stated as a conjecture supported by two
exact moments; we do not claim it as a theorem.
\end{abstract}

\noindent\textbf{MSC 2020.} Primary 05C80, 05C07; secondary 60C05, 05C35.\\
\textbf{Keywords.} random $k$-out graph, subgraph counts, triangles, clustering,
stars, orientation invariant, exact formulae.

\section{Introduction}
\subsection{The model and motivation}

The \emph{random $k$-out graph} $G(n,k\text{-out})$ (in the symmetrized form used
throughout the literature, e.g.\
\cite[Ch.~22]{friezeKaronski2015}, \cite{soodYagan2020}) is constructed as follows.
Each of the $n$ vertices $v$ independently and uniformly selects a $k$-subset
$S_v\subseteq [n]\setminus\{v\}$; vertices $u$ and $v$ are adjacent if either
$u\in S_v$ or $v\in S_u$. The resulting simple graph has exactly between $nk/2$ and
$nk$ edges ($\mathbb E[\# \text{edges}] = nk - \tfrac{nk^2}{2(n-1)}$ exactly,
Proposition~\ref{prop:edges}) — a degree-bounded, connected-with-high-probability
($k\ge 2$) sparse model.

$G(n,k\text{-out})$ is a workhorse model in situations where each node has a hard
budget of $\approx k$ directed links: random pairwise key predistribution in sensor
networks \cite{eschenauerGligor2002,yaganMakowski2009,soodYagan2020}, payment channel
networks, and related designs; see \cite{elumarSoodYagan2023} and the references in
\cite{soodYagan2020}. Its \emph{global} behaviour is well understood: connectivity is
achieved almost surely for every $k\ge2$
\cite{fennerFrieze1982,friezeKaronski2015}, with sharp finite-$n$ bounds
\cite{soodYagan2020}; perfect matchings and Hamilton cycles occur
with high probability \cite[Ch.~22]{friezeKaronski2015}; and $k$-out
subgraphs of large hosts are $k$-connected w.h.p.\
\cite{friezeJohansson2014}.

\subsection{The gap: local structure}
By contrast, the \emph{local} structure — the expected number of copies of a fixed
small graph $H$ (triangle, $4$-cycle, star, path, $\dots$) at finite $(n,k)$ — has
not, to our knowledge, been developed. The chapter of \cite{friezeKaronski2015}
devoted to $k$-out graphs (Ch.~22) contains connectivity, perfect matchings,
Hamilton cycles, and nearest-neighbour graphs, and no counts of fixed subgraphs
(checked against the authors' full text; Section~\ref{sec:related}). The recent
triangle analyses for other correlated models — random key graphs
\cite{yaganMakowski2009,yaganMakowski2017}, clustered random graphs
\cite{newman2009}, and the clustering-index estimates of
\cite{islakYesiloglu2024} (which treat Erd\H os--R\'enyi expressions and simulate
regular/Barab\'asi--Albert/Watts--Strogatz models, without the $k$-out model) — do
not cover this model either.

The question is non-trivial and the answer is not a transposition of the
$G(n,p)$ computation, because of a genuine structural feature: the maker choices
\emph{within} one chooser are negatively correlated,
\begin{equation}\label{eq:negcorr}
\Pr\bigl(v\in S_u \wedge w\in S_u\bigr)=\frac{k(k-1)}{(n-1)(n-2)}
\;<\;\Bigl(\frac{k}{n-1}\Bigr)^{2},
\end{equation}
so the per-vertex choices do not behave as independent coin flips, and a naive
$G(n,p)$-style ``independence'' computation can overshoot the truth by factors
of several (Table~\ref{tab:naive}: at $k=1$, $n=200$ the naive estimate is
$3.97\times$ the exact value).

\subsection{Contributions and overview}
We establish, for every fixed finite simple graph $H$ ($v$ vertices, $e$ edges) and
all finite $(n,k)$ with $1\le k\le n-1$:

\begin{enumerate}
\item \textbf{Exact general formula} (Theorem~\ref{thm:general}): an alternating
per-edge inclusion--exclusion whose summands are exact products of per-chooser
avoidance probabilities. It computes
$\mathbb E[\#\{\text{labelled copies of }H\}]$ exactly at any $(n, k)$.
\item \textbf{Triangles} (Theorems~\ref{thm:tri} and \ref{thm:variance}):
the exact closed form
$P(\triangle) = 1-3A^2+3A^2B-B^3$ (with $A=\frac{\binom{n-2}{k}}{\binom{n-1}{k}}$,
$B=\frac{\binom{n-3}{k}}{\binom{n-1}{k}}$), an alternate identity exposing the
negative correction to the naive value, exact series coefficients
$(4k^3-3k^2)/3,\; -\frac{k^2(k-1)(4k-3)}{2},
\frac{k^2(6k^3-39k^2+46k-15)}{6}$, exact $k=1$ specials
($P(\triangle)=2/(n-1)^3$), and the exact second moment of the triangle count with
$\operatorname{Var}(T)\to(4k^3-3k^2)/3$ and $1/n$-term
$-\frac{k^2(16k^2-21k+7)}{2}$.
\item \textbf{Matched-density Erd\H{o}s--R\'enyi contrast}
(Proposition~\ref{prop:er}): with the edge density $q=2p-p^{2}$,
$p=\tfrac{k}{n-1}$, matched between the two models, the ratio of expected triangle
counts tends to $1-\tfrac{3}{4k}$ (e.g.\ $\tfrac14$ at $k=1$,
$\tfrac78$ at $k=4$), with explicit $1/n$ corrections; the $4$-cycle ratio tends to
$1-\tfrac1k+\tfrac{1}{8k^2}$.
\item \textbf{Stars} (Theorem~\ref{thm:stars}): an exact closed form for stars
$K_{1,m}$ and the exact orders: the first non-zero term of the star probability at
fixed vertices is exactly of orders $n^{-m}$, so
$\mathbb E[\#\text{unl. }K_{1,m}] = \frac{c_m(k)}{m!}\,n(1+o(1))$ with explicit
polynomials $c_m(k)$ up to $m=6$ (Appendix, Table~\ref{tab:cm}) and, in general,
$c_m(k)=\sum_{j}\binom{m}{j}(k)_{j}\,k^{m-j}$.
\item \textbf{General asymptotics via an orientation invariant}
(Theorem~\ref{thm:omega}): for any fixed $H$,
\[
\mathbb E[\#\{\text{labelled copies of }H\}]\;/\; n^{v-e}\;\longrightarrow\;
\Omega_k(H)\;=\;\sum_{\sigma}\prod_{x}(k)_{t_x(\sigma)},
\]
the sum over the $2^{e}$ assignments $\sigma$ of each edge to exactly one of its
endpoints, with $t_x(\sigma)$ the number of edges assigned to $x$.
$\Omega_k(H)=0$ iff no assignment gives every vertex owned degree at most $k$; the
proof is a two-sided Bonferroni argument and is complete.
\item \textbf{Verification stack} (Section~\ref{sec:repro}): all closed forms
cross-checked symbolically over the rationals, by exhaustive enumeration over
\emph{all} choice profiles at small $n$ (agreement to machine zero, e.g.\
$\operatorname{Var}(T)=\tfrac{15985}{5832}$ at $(n,k)=(5,2)$), and by Monte Carlo at
large $n$ (Table~\ref{tab:mc}: worst $|z|=1.43$ on means, $1.41$ on variances).
\end{enumerate}

\emph{Honesty statement on novelty.} We are \emph{not} aware of any published
closed-form expression for the expected number of copies of a fixed graph
$H$ (with $e\ge2$) in the symmetrized random $k$-out graph, at finite $(n,k)$ or in
asymptotic form. This negative statement is the result of the targeted search
described in Section~\ref{sec:related}; we cannot exclude the existence of
unindexed or overlooked prior work, and we phrase our contribution accordingly.

\subsection{Notation}
Throughout, $n\ge 2$, $1\le k\le n-1$, and $(a)_m = a(a-1)\cdots(a-m+1)$ is the
falling factorial (with $(a)_0=1$), $\binom{r}{s}=0$ if $r<s$ or $s<0$. All labels
are of the template $H$'s vertices; a \emph{labelled copy} of $H$ in $G$ is an
injective map $\varphi:V(H)\to[n]$ such that $\varphi(u)\varphi(v)\in E(G)$ for every
$uv\in E(H)$; the number of \emph{unlabelled} copies sitting on a given vertex set is
the labelled count divided by $|\mathrm{Aut}(H)|$. The expectation is over the
random chooser sets $(S_v)_{v\in[n]}$, independent across vertices.

\section{Preliminaries}\label{sec:model}

\begin{definition*}[Model]
$G(n,k\text{-out})$: every $v\in[n]$ draws $S_v$ uniformly from the
$\binom{n-1}{k}$ $k$-subsets of $[n]\setminus\{v\}$; the chooser sets
$(S_v)$ are independent. An edge $uv$ exists iff $u\in S_v$ or $v\in S_u$.
\end{definition*}

For a fixed unordered pair $\{u,v\}$, $\Pr(uv\in E) = q := 2p-p^{2}$ with
$p=\tfrac{k}{n-1}$: either $v$ chose $u$ (probability $p$) or else $v$ did not choose
$u$ (prob. $1-p$) and $u$ chose $v$ (prob. $p$, independent across choosers).
Note $q=1-A^{2}$ with $A:=\binom{n-2}{k}/\binom{n-1}{k}=(1-\tfrac{k}{n-1})$.

\begin{lemma}[Within-chooser negative correlation]\label{lem:negcorr}
For fixed $u$ and distinct $v,w\ne u$,
\[
\Pr(v\in S_u \wedge w \in S_u) = \frac{\binom{n-3}{k-2}}{\binom{n-1}{k}}
=\frac{k(k-1)}{(n-1)(n-2)} < \Bigl(\frac{k}{n-1}\Bigr)^{2} = \Pr(v\in S_u)\Pr(w\in S_u).
\]
\end{lemma}
\begin{proof}
Hypergeometric: the count of $k$-subsets of $[n]\setminus\{u\}$ containing both
$v,w$ is $\binom{n-3}{k-2}$. The ratio expands to $\frac{k(k-1)}{(n-1)(n-2)}$, and
$\frac{k(k-1)}{(n-1)(n-2)}<\frac{k^2}{(n-1)^2}$ for $k\ge1$. \qed
\end{proof}

Chooser independence (across different vertices $u$) is exact; it is only the
\emph{within}-chooser dependencies (Lemma~\ref{lem:negcorr}) that break the
$G(n,p)$ analogy. Our exact formulae keep track of them through exact
avoidance probabilities.

\begin{proposition}[Edges]\label{prop:edges}
\[
\mathbb E[\,\#\text{edges}\,] = \binom{n}{2}\,q = nk-\frac{nk^2}{2(n-1)}\quad\text{exactly.}
\]
\end{proposition}

\section{The general subgraph-count formula}\label{sec:general}

Fix a template $H$ with vertex set \(\{1,\dots,v\}\) and edge list
$E(H)$. For a subset $F\subseteq E(H)$ let $d_x(F)$ be the degree of $x$ in the
spanning subgraph $(V(H),F)$.

\begin{theorem}[Exact expected labelled copy count]\label{thm:general}
For fixed $(n,k)$ and any fixed injective map $\varphi:\, V(H)\to [n]$,
\[
\Pr\bigl(\varphi \text{ is an embedding}\bigr)
\;=\;\sum_{F\subseteq E(H)}(-1)^{|F|}\ \prod_{x\in V(H)}
\frac{\binom{\,n-1-d_x(F)\,}{k}}{\binom{n-1}{k}},
\]
hence, exactly,
\[
\mathbb E\bigl[\,\#\{\text{labelled copies of }H\}\,\bigr]
= (n)_{v}\ \sum_{F\subseteq E(H)}(-1)^{|F|}\ \prod_{x}
\frac{\binom{\,n-1-d_x(F)\,}{k}}{\binom{n-1}{k}},
\]
where the convention $\binom{r}{k}=0$ for $r<k$ applies. For fixed $k$ and $H$, the
summands are rational numbers (exact), and the formula involves at most
$2^{e}$ terms.
\end{theorem}

\begin{proof}
By symmetry over injective maps, the probability for any fixed $\varphi$ equals the
probability for the identity embedding; write $x$ for the vertex of label
$\varphi(x)$. For $uv \in E(H)$ let $X_{uv}$ be the event that $uv$ is an edge of
$G(n,k\text{-out})$. Then
$\Pr(\forall uv\in E(H):\, X_{uv}) =\ \sum_{F \subseteq E(H)}(-1)^{|F|}\,
\Pr\bigl(\bigcap_{uv\in F}\overline{X_{uv}}\bigr)$
by inclusion--exclusion.

Now $\bigcap_{uv \in F}\overline{X_{uv}}$ holds iff every chooser $x$ avoids each of
its $d_x(F)$ template-neighbours. The chooser sets are independent across $x$, and
the probability that $x$'s uniformly random $k$-subset of $[n]\setminus\{x\}$
avoids a fixed set of $d_x(F)$ forbidden targets is exactly
$\binom{n-1-d_x(F)}{k}/\binom{n-1}{k}$ (hypergeometric, exact); multiplying this
over the (independent) choosers $x\in V(H)$ gives
$\prod_x \binom{\,n-1-d_x(F)\,}{k}/\binom{n-1}{k}^{\,v}$.
The expectation follows because for a uniform injective $\varphi$, the
$\varphi$-images of the $v$ template vertices are the entries of a uniformly random
$v$-tuple of distinct vertices, giving the factor $(n)_v$. \qed
\end{proof}

\begin{remark}
The formula has the shape of a $G(n,p)$-type inclusion--exclusion \emph{except}
that the ``per-edge exponentially-closed'' factors of $G(n,p)$
(factors $(1-p^2)^{|\cdot|}$) are replaced by \emph{per-chooser} products of exact
avoidance probabilities; the difference is precisely what survives of the negative
correlation of Lemma~\ref{lem:negcorr}. For later use we also record that all
predictions of Theorem~\ref{thm:general} were verified independent of it, by
exhaustive enumeration over all choice profiles
(Section~\ref{sec:repro}: every comparison has absolute difference $\le 2.2\cdot
10^{-15}$, mostly exactly $0$).
\end{remark}

\section{Triangles: exact formula, asymptotics, and the naive gap}\label{sec:tri}

Let $T$ be the number of (unlabelled) triangles: $T=\sum_{\{a,b,c\}}\mathbf 1\{abc
\text{ triangle}\}$.

\begin{theorem}[Exact triangle probability]\label{thm:tri}
Let
\[
A=\frac{\binom{n-2}{k}}{\binom{n-1}{k}} = 1-\frac{k}{n-1},\qquad
B=\frac{\binom{n-3}{k}}{\binom{n-1}{k}}
=\Bigl(1-\frac{k}{n-1}\Bigr)\Bigl(1-\frac{k}{n-2}\Bigr).
\]
Write $C:=\tfrac{n-2-k}{n-2}$, so that $B=AC$. Then exactly
\begin{equation}\label{eq:triflat}
P(\triangle):=\Pr(\text{fixed triple }\{a,b,c\}\text{ spans a triangle})
=1-3A^{2}+3A^{2}B-B^{3},
\end{equation}
and, equivalently,
\begin{equation}\label{eq:triident}
P(\triangle)\;=\;\underbrace{(1-A^{2})^{3}}_{\text{naive }G(n,p)\text{ value}}\;+\;
A^{3}\,(C-A)\,\bigl(3-A^{2}-AC-C^{2}\bigr),
\end{equation}
where the second term is $\le 0$ (since $C<A$), a \emph{negative correction} to the
naive value.
\end{theorem}

\begin{proof}
By Theorem~\ref{thm:general} specialized to the triangle (edges $ab,bc,ca$), a subset
$F\subseteq\{ab,bc,ca\}$ of size $0,1,2,3$ yields $(d_a,d_b,d_c)$ equal to
$(0,0,0)$; $(1,1,0)$-permutations; $(2,1,1)$-permutations; and $(2,2,2)$,
respectively. There are $1$, $3$, $3$, $1$ such subsets, hence
$P(\triangle)=1-3A^2+3A^2B-B^3$, since the middle product is
$\frac{\binom{n-3}{k}}{\binom{n-1}{k}}\cdot
\frac{\binom{n-2}{k}}{\binom{n-1}{k}}\cdot\frac{\binom{n-2}{k}}{\binom{n-1}{k}} = B\,A^2$.
The identity \eqref{eq:triident} is an algebraic rearrangement using $B=AC$
and $q=1-A^2$; both the relation $B\equiv AC$ and the equivalence
\eqref{eq:triident}$\equiv$\eqref{eq:triflat} are machine-checked symbolically.
The sign of the correction: $C - A = \tfrac{n-2-k}{n-2} - \tfrac{n-1-k}{n-1}
= -\tfrac{k}{(n-1)(n-2)}< 0$, while $3-A^2-AC-C^2 >0$ for the relevant range
(it is $3 - (A^2+AC+C^2)$ with $A,C \in (0,1)$, so $A^2+AC+C^2 < 3$).
\qed
\end{proof}

\begin{theorem}[Asymptotics of the mean]\label{thm:triasy}
As $n\to\infty$ with $k$ fixed, with $T$ the triangle count,
\[
\mathbb E T \;=\;e_0+\frac{e_1}{n}+\frac{e_2}{n^{2}}
+O\!\left(\frac{1}{n^{3}}\right),
\]
where
\[
e_0=\frac{4k^{3}-3k^{2}}{3},\qquad
e_1=-\frac{k^{2}(k-1)(4k-3)}{2},\qquad
e_2=\frac{k^{2}(6k^{3}-39k^{2}+46k-15)}{6},
\]
and the next coefficient is
$\ e_3=-\frac{k^{2}(k^{4}-27k^{3}+99k^{2}-96k+27)}{6}$
(by which the truncated series at $1/n^{2}$ is controlled).
In particular \ $\lim \mathbb E T = (4k^{3}-3k^{2})/3$, i.e.\ $\Theta(1)$ for fixed
$k$: the mean triangle count does not grow with $n$ at all.
\end{theorem}

The coefficients were derived by exact series expansion of
$\binom{n}{3}P(\triangle)$ (Theorem~\ref{thm:tri}) in $1/n$ and are machine-checked
symbolically; the truncation at $e_2$ was additionally checked against exact
Fraction evaluations at $n=400$ for $k=1,\dots,5$, with residuals bounded by
$|e_3|/n^3$ (Section~\ref{sec:repro}).

\begin{corollary}[Exact $k=1$ specials]\label{cor:k1}
For $k=1$:
\[
P(\triangle) = \frac{2}{(n-1)^{3}}\quad\text{and}\quad
P(\text{4-cycle}) = \frac{2}{(n-1)^{4}}
\quad\text{for all } n.
\]
\end{corollary}
These were simplified to closed form from the general formulae and
machine-checked symbolically.

\begin{remark}[Exact form of the naive gap]\label{rem:gap}
The naive $G(n,p)$-style prediction for a triple is $q^{3}=(1-A^{2})^{3}$; the exact
gap is, by identity~\eqref{eq:triident},
\[
q^{3}-P(\triangle) = -A^{3}(C-A)\bigl(3-A^{2}-AC-C^{2}\bigr)\;\ge\;0,
\]
a rational function of $(n,k)$ that starts contributing only from order $n^{-5}$
(the orders $n^{-3},n^{-4}$ coefficients of the gap are $0$);
its first non-zero terms are $\ \frac{6k^{2}}{n^{5}}
-\frac{3k^{2}(7k-9)}{n^{6}}+\cdots$. Nevertheless the \emph{relative} error of the
naive prediction is large because the truth itself is small; see
Table~\ref{tab:naive}.

\medskip
At matched edge density $q$ (same expected edges as $G(n,p)$), the exact
$1/n$-expansion of the ratio is
\[
\frac{P(\triangle)}{q^{3}}
= \underbrace{1-\frac{3}{4k}}_{r_0}
+ \underbrace{\frac{3(4k-3)}{8k}}_{r_1}\frac{1}{n}+O(n^{-2}),
\]
which is analysed next (Proposition~\ref{prop:er}).
\end{remark}

\begin{table}[t]
\centering
\caption{Exact $\mathbb E T$ vs.\ the naive $(G(n,p)$-style) prediction
$(1-A^{2})^{3}\binom{n}{3}$ at $n=200$. From \texttt{results/expectations\_table.csv}.}
\label{tab:naive}
\begin{tabular}{rrrrr}
\toprule
$k$ & exact $\mathbb ET$ & naive & rel.\ error & naive/exact\\
\midrule
1 & 0.333325 & 1.323275 & 2.9699 & 3.97\\
2 & 6.616147 & 10.506404 & 0.5880 & 1.59\\
3 & 26.592529 & 35.191162 & 0.3233 & 1.32\\
4 & 67.768683 & 82.784150 & 0.2216 & 1.22\\
5 & 137.416009 & 160.459789 & 0.1677 & 1.17\\
6 & 242.574103 & 275.163298 & 0.1343 & 1.13\\
\bottomrule
\end{tabular}
\end{table}

\section{The Erd\H os--R\'enyi contrast at matched density}\label{sec:er}

The natural comparison uses the same expected edge density
$q=2p-p^{2}$, $p=\tfrac{k}{n-1}$, in the Erd\H os--R\'enyi model
$G(n,q)$, for which $\mathbb E T_{ER} = \binom{n}{3}q^{3}$ exactly.

\begin{proposition}[Triangle and 4-cycle ratios]\label{prop:er}
As $n\to\infty$ with $k\ge 1$ fixed:
\[
\rho_\triangle:=\frac{\mathbb E T_{k\text{-out}}}{\mathbb E T_{ER}}
= \frac{P(\triangle)}{q^{3}}
\;\longrightarrow\; 1-\frac{3}{4k},
\]
with $1/n$-term $r_1=\frac{3(4k-3)}{8k}$ and next term
$-\frac{3(3k^{2}-21k+14)}{16k}\cdot n^{-2}$; and
\[
\rho_{C_4}:=\frac{\mathbb E[\#C_4]_{k\text{-out}}}{\mathbb E[\#C_4]_{ER}}
\;\longrightarrow\;1-\frac{1}{k}+\frac{1}{8k^{2}},
\]
with $1/n$-term $\frac{8k^{2}-8k+1}{4k^{2}}$.
\end{proposition}

\begin{proof}
$P(\triangle)/q^{3}$ and the closed form for $P(C_4)$
(computed in Lemma~\ref{lem:c4} below) divided by $q^{4}$ are rational functions of
$(n,k)$. The expansions are exact series expansions of those two ratios in $1/n$
(machine-checked symbolically; the leading terms simplify to the stated
polynomial expressions in $k$).\qed
\end{proof}

\begin{lemma}[4-cycle closed form]\label{lem:c4}
\[
P(C_4)=1-4A^{2}+4A^{2}B+2A^{4}-4A^{2}B^{2}+B^{4},
\]
for a fixed labelled 4-cycle $\{a b c d\}$, with
$\mathbb E[\#\text{labelled }C_4] = (n)_4 P(C_4)$ and
$\ \mathbb E[\#\text{unlabelled }C_4] = 3\binom{n}{4}P(C_4)$.
\end{lemma}
\begin{proof}
Apply Theorem~\ref{thm:general} to the template $C_4$; the signed sum over the
$2^{4}=16$ subsets groups into classes by edge pattern (verified symbolically: the
Theorem~\ref{thm:general} sum equals the closed form as an identity in
$(n,k)$).\qed
\end{proof}

\begin{table}[t]
\centering
\caption{Triangles: exact series vs.\ limits ($n=1600$ row of
\texttt{results/asymptotic\_convergence.csv}).}
\label{tab:asy}
\begin{tabular}{rrrrrr}
\toprule
$k$ & $\mathbb E_{n=1600}$ & $e_0$ & $(\mathbb E-e_0)\,n$ & $e_1$ & $\rho_\triangle$ ($\to 1-\tfrac3{4k}$)\\
\midrule
1 & 0.3333332 & $\tfrac13$ & $-2.1\cdot10^{-4}$ & $0$ & 0.250235 ($0.250$)\\
2 & 6.6604086 & $\tfrac{20}{3}$ & $-10.0129$ & $-10.0000$ & 0.625587 ($0.625$)\\
3 & 26.9493363 & $27$ & $-81.0619$ & $-81.0000$ & 0.750704 ($0.750$)\\
4 & 69.1382595 & $\tfrac{208}{3}$ & $-312.1181$ & $-312.0000$ & 0.813262 ($0.8125$)\\
5 & 141.1354011 & $\tfrac{425}{3}$ & $-850.0249$ & $-850.0000$ & 0.850797 ($0.850$)\\
\bottomrule
\end{tabular}
\end{table}

\subsection{Interpretation}
For every fixed $k\ge1$, the $k$-out graph has \emph{asymptotically}
$\left(1-\tfrac{3}{4k}\right)$-fraction of the triangles that an
Erd\H{o}s--R\'enyi graph of equal expected edge density has: $\tfrac14$ at $k=1$,
$\approx\tfrac23$ approaching as $k\to\infty$. This is a large, real effect in the
opposite direction from most ``good clustering'' enhancers (Barab\'asi--Albert,
Watts--Strogatz): the $k$-out budget spreads links thin and thereby makes the
neighbourhoods of vertices locally tree-like, and the effect is \emph{stronger
for small $k$}. Together with the exact finite-$n$ formulae, the constants
$1-\tfrac{3}{4k}$ and $1-\tfrac{1}{k}+\tfrac{1}{8k^{2}}$ quantify what practitioners
trade away (or gain, for triangle-suppression purposes) when replacing an
Erd\H{o}s--R\'enyi null by a degree-bounded $k$-out null.

\section{Stars: exact cancellation and explicit constants}\label{sec:stars}

We use $K_{1,m}$ (a star with $m$ leaves) at \emph{fixed} vertices
$\{u$ (center), $v_1,\dots,v_m\ ($leaves$)\}$ and $N_{K_{1,m}}$ for the count of
unlabelled copies.

\begin{theorem}[Stars]\label{thm:stars}
Fix $m\ge2$ and $k\ge1$. Exactly,
\begin{equation}\label{eq:starclosed}
P(K_{1,m})
=\sum_{j=0}^{m}(-1)^{j}\binom{m}{j}\,\frac{\binom{n-1-j}{k}}{\binom{n-1}{k}}\,
\Bigl(\frac{\binom{n-2}{k}}{\binom{n-1}{k}}\Bigr)^{j},
\end{equation}
and the first non-zero term of $P(K_{1,m})$ in powers of $1/n$ is \emph{exactly of
order} $n^{-m}$, with finite non-zero constant
\begin{equation}\label{eq:cmconst}
\lim_{n\to\infty} n^{m} P(K_{1,m}) \;=\; c_m(k)
\;:=\;\sum_{j=0}^{m}\binom{m}{j}\,(k)_{j}\,k^{m-j} \;>\;0.
\end{equation}
Consequently
\[
\mathbb E\bigl[\,N_{K_{1,m}}\bigr]=\frac{(n)_{m+1}}{m!}\,P(K_{1,m})
=\frac{c_m(k)}{m!}\,n\,\bigl(1+o(1)\bigr) \quad\text{as }n\to\infty,
\]
with explicit polynomial constants $c_2(k)=k(4k-1)$,
$c_3(k)=2k(4k^2-3k+1)$, $c_4(k)=k(16k^3-24k^2+19k-6)$, $c_5(k)=
2k(16k^4-40k^3+55k^2-40k+12)$, $c_6(k)=k(64k^5-240k^4+500k^3-615k^2+418k-120)$
(Table~\ref{tab:cm}).
\end{theorem}

\begin{proof}[Proof]
\eqref{eq:starclosed} is the specialization of Theorem~\ref{thm:general} to the
star: a subset $F$ of the $m$ star edges consists of $j$ edges, all incident to the
center; the spanning degrees are $d_{\text{center}} = j$ and $d_{v_i}=1$ for the
$j$ selected leaves and $0$ for the others. Hence each $F$ of size $j$ contributes
sign $(-1)^{j}$ times the product
$\frac{\binom{n-1-j}{k}}{\binom{n-1}{k}}\,
\bigl(\frac{\binom{n-2}{k}}{\binom{n-1}{k}}\bigr)^{j}$, and there are
$\binom{m}{j}$ such subsets, summing to \eqref{eq:starclosed}. The identity between
\eqref{eq:starclosed} and the general formula is machine-checked for
$m=2,\dots,6$ (as identities in $(n,k)$ over the rationals).

For \eqref{eq:cmconst}: the orientation-invariant Theorem~\ref{thm:omega} below
(stated and proved in Section~\ref{sec:omega}) applies with $H=K_{1,m}$,
$v=m+1$, $e=m$, so $v-e=1$:
$\mathbb E[\#\text{labelled copies}]=\Omega_k(K_{1,m})\,n\,(1+o(1))$ with
$\Omega_k(K_{1,m})=\sum_{j}\binom{m}{j}(k)_{j}k^{m-j}$ (each edge is assigned to
either its center or its leaf; assignments with $j$ center-owned edges contribute
$\binom{m}{j}(k)_{j}k^{m-j}$). Dividing by $|\mathrm{Aut}(K_{1,m})|=m!$ and using
$\mathbb E[N_{K_{1,m}}]=(n)_{m+1}P(K_{1,m})/m!$ yields
$n^{m}P(K_{1,m})\to c_m(k)$, which is finite and positive: all summands of
\eqref{eq:cmconst} are non-negative and the $j=0$ summand is $k^{m}>0$.
\end{proof}

\begin{remark}
We emphasize two points of honesty about Theorem~\ref{thm:stars}.
(a) The identity \eqref{eq:cmconst}, with the closed-form constants verified
numerator-by-numerator in \texttt{src/experiments.py} (\texttt{STAR\_C},
$m\le6$), rests on Theorem~\ref{thm:omega}, whose proof is complete in this paper
(Section~\ref{sec:omega}).
(b) The exact orders claim — that \emph{no} term of order between $n^{0}$ and
$n^{-m}$ survives in ~\eqref{eq:starclosed} — \emph{follows} from
\eqref{eq:cmconst} with $c_m(k)>0$ finite: any lower-order non-zero term would blow
up $n^{m}P(K_{1,m})$. So the exact-order statement is \emph{equivalent} to the
finite-non-zero-limit statement, and both carry the weight of
Theorem~\ref{thm:omega}. The $m\leq6$ cases were in addition directly
series-expanded symbolically, and the constants $c_m(k)$ assessed numerically
with the residual decaying like $1/n$ (Table~\ref{tab:cm}; e.g.\
$n^{6}P(K_{1,6})$ at $k=2$: $930.1018$ at $n=1600$ vs.\ $c_6(2)=928$).
\end{remark}

\begin{table}[t]
\centering
\caption{Stars: $n^{m}P(K_{1,m})$ converging to $c_m(k)$ ($n=1600$ rows of
\texttt{results/star\_cancellation.csv}; residuals $O(1/n)$).}
\label{tab:cm}
\begin{tabular}{rrrrr}
\toprule
$m$ & $k$ & $n^{m}P$ at $n=1600$ & $c_m(k)$ & $\mathbb E[N_{K_{1,m}}]$ at $n=1600$\\
\midrule
2 & 2 & 14.0037 & 14 & 11182.0\\
2 & 3 & 32.9887 & 33 & 26341.5\\
3 & 2 & 44.0300 & 44 & 11697.3\\
3 & 3 & 167.9715 & 168 & 44624.6\\
4 & 2 & 128.1500 & 128 & 8490.1\\
4 & 3 & 801.1104 & 801 & 53074.3\\
5 & 2 & 352.6002 & 352 & 4657.4\\
5 & 3 & 3619.9307 & 3618 & 47814.9\\
6 & 2 & 930.1018 & 928 & 2039.9\\
6 & 3 & 15648.4509 & 15633 & 34320.3\\
\bottomrule
\end{tabular}
\end{table}

\section{The orientation invariant: general asymptotics}\label{sec:omega}

Fix a template $H$ with $v$ vertices and $e$ edges. An \emph{edge-owner
assignment} is a function $\sigma: E(H)\to V(H)$ with $\sigma(uv)\in\{u,v\}$ for
every edge $uv$; let $t_x(\sigma) = |\sigma^{-1}(x)|$ be the number of edges assigned
to $x$. Assignments correspond bijectively to \emph{orientations} of $H$'s edges
(towards the fixed endpoint ``owner'', or equivalently edge $\to$ one of its two
endpoints). Write $(k)_t$ for the falling factorial.

\begin{theorem}[Orientation invariant]\label{thm:omega}
For every fixed finite simple $H$ and fixed $k\ge1$,
\[
\mathbb E\bigl[\,\#\{\text{labelled copies of }H\}\,\bigr]
=\Omega_k(H)\,n^{v-e}\,(1+o(1))
\quad\text{ as } n\to\infty,
\]
where
\[
\Omega_k(H)\;=\;\sum_{\sigma\in\{0,1\}^{E(H)}}\
\prod_{x\in V(H)}(k)_{t_x(\sigma)}.
\]
Moreover, $\Omega_k(H)=0$ if and only if \emph{no} assignment $\sigma$ satisfies
$t_x(\sigma)\le k$ for all $x$; equivalently, iff every orientation of $H$ leaves
some vertex owning $\ge k+1$ edges; equivalently (via the orientation identity) iff
$H$ cannot be oriented such that each vertex has out-degree at most $k$ —
i.e.\ $H$ admits no orientation of maximum owned-degree $\le k$.
\end{theorem}

\begin{proof}
Fix the template labels (the identity embedding, by symmetry). For an assignment
$\sigma$ let $Y_\sigma$ be the event that, for every edge $uv \in E(H)$, the
\emph{owner} $\sigma(uv)$ has the \emph{other} endpoint of $uv$ in its chosen set:
if $uv$ is owned by $u$, this means $v\in S_u$. If $Y_\sigma$ holds, then every
template edge is present (each edge is chosen by its owner), so
$Y_\sigma\subseteq\{\text{all template edges present}\}$; conversely, if all edges
are present, then every edge has at least one choosing endpoint, and choosing one
such endpoint per edge extends to an assignment; hence
$\{\text{all template edges present}\}=\bigcup_\sigma Y_\sigma$.

The chooser sets are independent across vertices, and the probability that a
particular chooser $x$ places at least its $t_x(\sigma)$ prescribed target-neighbours
into its $k$-subset is exactly $(k)_{t_x(\sigma)}/(n-1)_{t_x(\sigma)}$
(the count of $k$-subsets containing $t$ specific targets divided by
$\binom{n-1}{k}$). Hence, by independence,
\[
\Pr(Y_\sigma) \;=\;\prod_{x}\frac{(k)_{t_x(\sigma)}}{(n-1)_{t_x(\sigma)}}\;.
\]
Since $\sum_x t_x(\sigma) = e$, we get
\[
(n)_{v}\,\Pr(Y_\sigma) =
\frac{(n)_{v}}{\prod_x (n-1)_{t_x(\sigma)}}\ \prod_x (k)_{t_x(\sigma)}
= n^{v-e}\,\Omega(\sigma)\,\bigl(1+O(n^{-1})\bigr),
\]
uniformly in $\sigma$, where $\Omega(\sigma)=\prod_x (k)_{t_x(\sigma)}$.
Summing over the $2^{e}$ assignments:
\[
(n)_{v}\,\Pr\Bigl(\bigcup_\sigma Y_\sigma\Bigr) = \Omega_k(H)\,n^{v-e}\,\bigl(1+O(n^{-1})\bigr)
\;\;+\;\; \text{(over-count correction)}.
\]
Two-sided control of the correction: for $\sigma\ne\tau$, the event
$Y_\sigma\cap Y_\tau$ requires at least one edge $uv$ to be chosen by both
endpoints (the edge on which they disagree must lie in both), an event of
probability $\frac{k(k-1)}{(n-1)(n-2)}=O(n^{-2})$; the remaining $e-1$ or fewer
edges contribute witnesses at $\ge n^{-1}$ each, so
$\Pr(Y_\sigma\cap Y_\tau) = O(n^{-e-1})$.
There are at most $\binom{2^{e}}{2} = O(1)$ pairs $(\sigma,\tau)$, hence
by the second Bonferroni inequality (applied to $\bigcup_\sigma Y_\sigma$),
\[
\Bigl|\ \Pr\Bigl(\bigcup_\sigma Y_\sigma\Bigr) - \sum_{\sigma}\Pr(Y_\sigma)\Bigr|
\ \le\ \sum_{\sigma<\tau}\Pr\bigl(Y_\sigma\cap Y_\tau\bigr)
= O(n^{-e-1}),
\]
so \ $(n)_v\Pr(\cup Y_\sigma) = \Omega_k(H)\, n^{v-e} (1+o(1)) + O(n^{v-e-1})$,
which is the claim.

The ``iff'' statement: $\Omega_k(H)=\sum_\sigma\prod_x(k)_{t_x(\sigma)}$ is a sum
of non-negative terms. If some assignment $\sigma$ satisfies $t_x(\sigma)\le k$
for every $x$, then all its factors $(k)_{t_x(\sigma)}$ are positive integers
(each at least $1$), so its summand is $\ge 1$ and $\Omega_k(H)\ge 1>0$.
Conversely, if every $\sigma$ has some $t_x(\sigma)\ge k+1$, then that factor
$(k)_{t_x(\sigma)}=0$ (the falling factorial vanishes once its length exceeds
$k$), so every summand — and hence $\Omega_k(H)$ — is $0$. \qed
\end{proof}

\begin{corollary}[Instances]\label{cor:omega}
Machine-verified instances (Section~\ref{sec:repro}):
\[
\Omega_k(K_3) = 8k^{3}-6k^{2},\qquad
\Omega_k(C_4) = 2k^{2}(8k^{2}-8k+1),\qquad
\Omega_k(K_{1,m}) = \sum_{j}\binom{m}{j}(k)_{j}k^{m-j},
\]
and, as a boundary example, $\Omega_1(K_4-e)=0$ (no orientation of the 5-edge
butterfly keeps every vertex's owned degree $\le1$), consistent with
$\Omega_1(K_3)=2>0$. Dividing by $|\mathrm{Aut}(H)|$ gives the unlabelled-limit;
e.g.\ $\lim \mathbb E T_{K_3}^{unlabelled} = (8k^3-6k^2)/6=(4k^3-3k^2)/3$ as in
Theorem~\ref{thm:triasy}.
\end{corollary}

\begin{remark}
(a) The orientation invariant makes the star constants transparent
($\Omega_k(K_{1,m}) = c_m(k)$, Theorem~\ref{thm:stars}). It also illustrates the
budget-jam threshold: $\Omega_k(H)$ contrasts sharply with $G(n,p)$, where the
leading behaviour for fixed $H$ at matched density is $q^{e}n^{v-e}$, i.e.\ an
$(2k/n)^{e}$-type scale: the $k$-out analogue $\Omega_k(H)$ \emph{cannot} grow
beyond what the per-vertex budget $k$ can support in owned-degree terms, which is a
strictly non-ER criterion — e.g.\ at $k=1$, $K_4$ is impossible ($\Omega_1(K_4)=0$),
while $K_3$ has $\Omega_1(K_3)=2$ copies-of-kind asymptotics.
(b) When $k \geq \Delta(H)$ (max degree), the identity-assignment ($\sigma = $ all
edges owned by their ``first'' endpoint under some orientation of $H$ where each
vertex owns $\le \Delta(H)\le k$) works, so $\Omega_k(H)>0$ always for fixed
$H$ and $k\ge \Delta(H)$.
\end{remark}

\section{Exact second moment of the triangle count}\label{sec:variance}

Let $T$ again be the (unlabelled) triangle count.

\begin{theorem}[Exact first two moments]\label{thm:variance}
For all $(n,k)$, exactly
\[
\mathbb E T^{2} \;=\; \binom{n}{3}P_3 \;+\; 12\binom{n}{4}P_{K_4-e} \;+\;
n\binom{n-1}{2}\binom{n-3}{2}P_{\text{bowtie}}\;+\;
\binom{n}{3}\binom{n-3}{3}P_{2\triangle},
\]
where $P_3 = P(\triangle)$ (Thm.~\ref{thm:tri}) and $P_{K_4-e}$, $P_{\text{bowtie}}$,
$P_{2\triangle}$ are the Theorem~\ref{thm:general} probabilities of the three
\emph{union graphs}: two triangles sharing an edge ($K_4$ minus an edge), sharing a
vertex (the bowtie), and disjoint. In particular $\operatorname{Var}(T) = \mathbb
E T^{2} - (\mathbb ET)^{2}$ is an exact rational function of $(n,k)$, and as
$n\to\infty$ with fixed $k$:
\[
\operatorname{Var}(T) = v_0+\frac{v_1}{n}+O(n^{-2}),\qquad
v_0=\frac{4k^{3}-3k^{2}}{3},\quad
v_1=-\frac{k^{2}(16k^{2}-21k+7)}{2},
\]
so that $\operatorname{Var}(T)\to \mathbb E T$ (Var $=$ mean asymptotically) with
negative $1/n$-correction $v_1 <0$ (underdispersion at every finite $n$ staged,
Table~\ref{tab:var}).
\end{theorem}

\begin{proof}
$T$ counts unordered triples; writing $T=\sum_{s}\mathbf 1_s$ over the
$\binom{n}{3}$ triples, we get
$\mathbb ET^{2} = \sum_{s,t}\Pr(E(s)\cup E(t))$, the (ordered) pair classes
handled by the \emph{overlap-pattern decomposition} of the two triples:
$s=t$ ($\binom{n}{3}$ pairs, union = one triangle);
$|s\cap t|=2$ ($12\binom{n}{4}$ ordered pairs; union $=K_4-e$);
$|s\cap t|=1$ ($n\binom{n-1}{2}\binom{n-3}{2}$ pairs; union $=$ bowtie);
$|s\cap t|=0$ ($\binom{n}{3}\binom{n-3}{3}$ pairs; union $=$ two disjoint
triangles). Each union's edge-present probability is a
Theorem~\ref{thm:general} value at the respective template sizes
($3$,$5$,$6$,$6$ edges on $3$,$4$,$5$,$6$ vertices). The series coefficients were
derived and checked symbolically (all identities machine-verified; the
$\operatorname{Var}$-series was additionally checked against exact Fraction
arithmetic at $n=400$).\qed
\end{proof}

\begin{conjecture}[Poisson-type limit]\label{conj:poisson}
For every fixed $k\ge1$, as $n\to\infty$, the triangle count
$T$ converges in distribution to $\text{Poisson}\bigl(\tfrac{4k^3-3k^2}{3}\bigr)$.
\end{conjecture}
We \underline{do not} prove Conjecture~\ref{conj:poisson}; we establish only the two
exact moments (Theorem~\ref{thm:variance}), whose agreement with the Poisson
moments (mean $=$ variance asymptotically) is suggestive, and we checked by Monte
Carlo that the exact mean and exact variance agree with the sampled moments in
Table~\ref{tab:mc} to within $|z|\le 1.43$ in all seven staged $(n,k)$ cells
(dispersion $\operatorname{Var}/\mathbb E$ at $n=200$: $0.985$ at $k=1$, $0.966$
at $k=2$, $0.947$ at $k=3$). A proof presumably requires a third-moment method or
the $G(n,p)$-style Janson machinery under the $k$-out correlated conditioning, and
is left open.

\begin{table}[t]
\centering
\caption{Monte Carlo vs.\ exact moments (\texttt{results/monte\_carlo.csv}; the
sampling corrects for the symmetrization via an OR-merge before the matrix
product, cross-checked against a naive Fisher-consistency sampler).}
\label{tab:mc}
\begin{tabular}{rrrrrrrr}
\toprule
$n$ & $k$ & reps & MC mean & exact $\mathbb E$ & $z$ & MC var / exact var & $z$\\
\midrule
50 & 2 & 20000 & 6.46805 & 6.45824 & $+0.58$ & 5.679 / 5.687 & $-0.13$\\
50 & 3 & 20000 & 25.32875 & 25.34068 & $-0.36$ & 21.744 / 21.437 & $+1.41$\\
100 & 2 & 20000 & 6.57570 & 6.56458 & $+0.63$ & 6.154 / 6.133 & $+0.34$\\
100 & 3 & 20000 & 26.21760 & 26.18013 & $+1.09$ & 23.662 / 23.669 & $-0.03$\\
200 & 3 & 20000 & 26.64355 & 26.59253 & $+1.43$ & 25.333 / 25.182 & $+0.60$\\
200 & 4 & 20000 & 67.80790 & 67.76868 & $+0.70$ & 63.251 / 63.114 & $+0.22$\\
500 & 3 & 4000 & 26.80175 & 26.83760 & $-0.44$ & 26.404 / 26.234 & $+0.29$\\
\bottomrule
\end{tabular}
\end{table}

\begin{table}[t]
\centering
\caption{Dispersion of the triangle count (\texttt{results/variance\_table.csv};
$k=1$,$2$ shown; $k=3$ staged in the file).}
\label{tab:var}
\begin{tabular}{rrrrr}
\toprule
$k$ & $n$ & $\operatorname{Var}/\mathbb E$ & $(\operatorname{Var}-v_0)n$ & $v_1$\\
\midrule
1 & 200 & 0.98515 & $-0.9916$ & $-1.0000$\\
2 & 200 & 0.96558 & $-55.6546$ & $-58.0000$\\
1 & 1000 & 0.99701 & $-0.9983$ & $-1.0000$\\
2 & 1000 & 0.99286 & $-57.5259$ & $-58.0000$\\
\bottomrule
\end{tabular}
\end{table}

\section{Verification and reproducibility}\label{sec:repro}

All numeric claims in this paper are reproduced by the public repository (see
below for the placement), which contains the source code, the generated data, and
the verification documents. The repository's verification stack:

\begin{enumerate}
\item \textbf{Symbolic}: \texttt{src/verify\_symbolic.py} — 36 machine-verified
identities over exact rational algebra (sympy), including Theorem~\ref{thm:general}
vs.\ closed forms for triangles/$C_4$/stars ($m\le6$), the identity
\eqref{eq:triident}, $B=AC$, the $k=1$ specials of Corollary~\ref{cor:k1}, the
series coefficients of Theorems~\ref{thm:triasy} and \ref{thm:variance}
(with $e_3$-based truncation error bounds), the edge identity, the matched-density
ratio constants of Proposition~\ref{prop:er}, and the orientation-invariant
instances of Corollary~\ref{cor:omega}.
\item \textbf{Exhaustive enumeration}: for small $n$, \emph{all}
$\prod_{v}\binom{n-1}{k}$ choice profiles are enumerated and the copy counts
matched exactly (exact-Fraction or float; 25 staged cases; largest difference
$2.2\cdot10^{-15}$; the $(n,k)=(5,2)$ second-moment case yields
$\operatorname{Var}(T) = \tfrac{15985}{5832}$ exactly on both sides).
\item \textbf{Monte Carlo}: at larger $n$, Fisher-consistent sampling
(symmetrized by an explicit OR-merge) with the staged cells of
Table~\ref{tab:mc}: worst $|z| = 1.43$ on means, $1.41$ on variances, over
4{,}000--20{,}000 replicates per cell.
\item \textbf{Unit tests}: \texttt{pytest} (\texttt{tests/test\_kout.py}),
43 tests, all passing (as of the release date recorded in the repository).
\end{enumerate}

\medskip
\noindent\textbf{Reproducibility statement.} The repository contains:
\texttt{src/} (formulae, closed forms, exact-Fraction paths, exhaustive
enumeration, Monte Carlo); \texttt{results/*.csv} (all tables behind
Figures/Tables in this paper); \texttt{research/} (the research process documents:
candidate scoring, literature review, evidence ledger, novelty assessment, hostile
review, limitations, journal selection); and \texttt{paper/} (this manuscript and
its bibliography). Running \texttt{python src/experiments.py} regenerates every
CSV from scratch (Monte Carlo seeded); running \texttt{python} with the script
\texttt{src/verify\_symbolic.py} re-runs the 36 symbolic checks;
\texttt{pytest} re-runs the 43 tests. Software: Python 3.14 with numpy, sympy (exact arithmetic), pytest.
No proprietary data are used; the literature checks cite publicly retrievable
sources with their access dates.

\section{Related literature, in more detail}\label{sec:related}

\paragraph{Global $k$-out structure.} Connectivity and Hamiltonicity of
$G(n,k\text{-out})$ and close relatives: \cite{fennerFrieze1982}
(connectivity of random $m$-orientable graphs and digraphs, Combinatorica 2(4),
1982); the textbook chapter \cite[Ch.~22]{friezeKaronski2015} (connectivity,
perfect matchings, Hamilton cycles, nearest-neighbour graphs; no subgraph-count
content — we verified the section structure against the authors' full text,
retrieved October 2026); finite-$n$ connectivity bounds \cite{soodYagan2020};
robustness and giant component \cite{elumarSoodYagan2023}; $k$-out subgraphs of
large hosts \cite{friezeJohansson2014}.

\paragraph{Triangle/clustering in correlated models.} For random key graphs
(Eschenauer--Gligor key pools), \ \cite{yaganMakowski2009} proved a zero-one law
for triangle presence and \cite{yaganMakowski2017} computed the clustering
coefficient and contrasted it with matched ER — the closest methodological
precedent to our Section~\ref{sec:er}, but on a different model (uniform key pools
produce a distinct dependence structure and a zero-one law, not our
$1-\tfrac{3}{4k}$ density-ratio effect). Clustered random graph constructions with
explicit triangle stubs: \cite{newman2009}; epidemic percolation on such
constructions: Sood's thesis \cite{sood2024thesis} (triangles as inserted building
blocks, ``doubly-Poisson'' pairwise joint distributions; the $k$-out model enters
only via connectivity remarks — checked against the full text, retrieved October
2026).
Clustering \emph{index} estimates are given in \cite{islakYesiloglu2024} for
Erd\H{o}s--R\'enyi with simulation for regular, Barab\'asi--Albert and
Watts--Strogatz models; the $k$-out model does not appear.

\paragraph{$G(n,p)$ baselines.} \ The expected numbers
$\binom{n}{v}p^{e}$, the Poisson/normal dichotomy of subgraph counts
\cite{jansonLuczakRucinski2000}, and the classical zero-one/threshold framework of
\cite{bollobas2001} are the backdrop for our matched-density comparison; none of
them applies verbatim because of Lemma~\ref{lem:negcorr}.

\paragraph{What is not known (to us).} Beyond the Poisson-type Conjecture
\ref{conj:poisson}: (a) exact higher moments or the full distribution of $T$;
(b) second moments for general $H$; (c) regimes with $k=k(n)\to\infty$ (dense
$k$-out), where new phenomena must appear since
$\mathbb E T\sim \tfrac{(4/3)k^3}{\cdot}$ would need to grow and cross the
$G(n,p)$ law; (d) clustering coefficients of the $k$-out model at fixed $k$ (which
are not the ratio $\rho_\triangle$ and involve degree conditioning).

\section{Limitations}\label{sec:limits}

(1) Our distributional results are first/second moment only;
Conjecture~\ref{conj:poisson} is not proven. (2) All asymptotic statements assume
fixed $k$ and fixed $H$; $k$ growing with $n$ is open. (3) The general-$m$ star
constants are explicit polynomials up to $m=6$ (tabulated); the general
formula $\sum_j\binom{m}{j}(k)_j k^{m-j}$ is proven via
Theorem~\ref{thm:omega} for all $m$. (4) The orientation-invariant proof is
complete, but its instance checks in the repository cover structured families
($K_3$, $C_4$, stars, $K_4-e$); a general-implementation cross-check for arbitrary
$H$ was not built in-session. (5) The novelty statement
(Section~\ref{sec:repro}) is a negative claim from a targeted search (documents in
the repository), not a proof of no-overlap: residual risks are recorded in the
evidence ledger.

\section*{Data and code availability}
All code, data (CSV), and research documentation are in the public repository
accompanying this work (see the reproducibility statement in
Section~\ref{sec:repro}).

\section*{Competing interests}
None. This work received no funding.

\appendix

\section{Star constants $c_m(k)$ for $m\le6$}\label{app:cm}
\[
\begin{aligned}
c_2(k) &= k\,(4k-1),\\
c_3(k) &= 2k\,(4k^{2}-3k+1),\\
c_4(k) &= k\,(16k^{3}-24k^{2}+19k-6),\\
c_5(k) &= 2k\,(16k^{4}-40k^{3}+55k^{2}-40k+12),\\
c_6(k) &= k\,(64k^{5}-240k^{4}+500k^{3}-615k^{2}+418k-120).
\end{aligned}
\]
Each admits the structural form
$c_m(k)=\sum_{j=0}^{m}\binom{m}{j}(k)_j\,k^{m-j}$
(Theorem~\ref{thm:omega}) and the tabulated closed forms agree with it exactly
(machine-checked for $m=2,\dots,6$). Values at $k=2,3$ (used in
Table~\ref{tab:cm}): $c_2: 14, 33$; $c_3: 44, 168$; $c_4: 128, 801$;
$c_5: 352, 3618$; $c_6: 928, 15633$.

\section{Verification protocol details}\label{app:verify}
The full verification documents (candidate scoring, literature review, evidence
ledger, novelty assessment, two hostile-review passes, limitations, journal
selection) accompany the repository under \texttt{research/}. The evidence ledger
records, per claim: the theorem number, the check that confirms it (symbolic /
enumeration / Monte Carlo / exact spot values), and any residual risks about the
literature survey. Two independent hostile passes over the manuscript were
recorded (pass 1: correctness-oriented attacks; pass 2: significance- and
reproducibility-oriented attacks), including the discovery and repair of a Monte
Carlo over-counting bug (documented; all Monte Carlo rows were regenerated after
the fix).

\begin{thebibliography}{12}

\bibitem{friezeKaronski2015}
A.~Frieze and M.~Karo\'nski,
\emph{Introduction to Random Graphs},
Cambridge University Press, Cambridge, 2015.
(Ch.~22 treats random $k$-out graphs: connectivity, perfect matchings, Hamilton
cycles; the authors' current full text is posted at
\url{https://www.math.cmu.edu/~af1p/BOOK.pdf}, accessed 2026-10-18.)

\bibitem{fennerFrieze1982}
T.~I.~Fenner and A.~M.~Frieze,
On the connectivity of random $m$-orientable graphs and digraphs,
\emph{Combinatorica} \textbf{2} (1982), no.~4, 347--359.

\bibitem{friezeJohansson2014}
A.~M.~Frieze and T.~Johansson,
On random $k$-out sub-graphs of large graphs,
arXiv:1405.2129 (2014).

\bibitem{bollobas2001}
B.~Bollob\'as,
\emph{Random Graphs}, 2nd ed.,
Cambridge University Press, Cambridge, 2001.

\bibitem{jansonLuczakRucinski2000}
S.~Janson, T.~{\L}uczak, and A.~Ruci\'nski,
\emph{Random Graphs},
Wiley-Interscience, New York, 2000.

\bibitem{islakYesiloglu2024}
{\"U}.~I{\c{s}}lak and B.~Ye{\c{s}}ilo{\u{g}}lu,
Analysis of clustering and degree index in random graphs and complex networks,
arXiv:2408.06451 (2024, v2 2025); published version in Journal of Applied
Probability, DOI 10.1017/jpr.2025.10063.

\bibitem{yaganMakowski2009}
O.~Ya{\u{g}}an and A.~M.~Makowski,
A zero-one law for the existence of triangles in random key graphs,
arXiv:0910.0499 (2009).

\bibitem{yaganMakowski2017}
O.~Ya{\u{g}}an and A.~M.~Makowski,
Counting triangles, tunable clustering and the small-world property in random key
graphs (extended version),
arXiv:1701.03758 (2017).

\bibitem{elumarSoodYagan2023}
E.~C.~Elumar, M.~Sood, and O.~Ya{\u{g}}an,
On the robustness, connectivity and giant component size of random $K$-out graphs,
arXiv:2311.02319 (2023).

\bibitem{soodYagan2020}
M.~Sood and O.~Ya{\u{g}}an,
Tight bounds for connectivity of random $K$-out graphs,
arXiv:2006.10638 (2020); an associated version appeared in IEEE ICC 2021,
DOI 10.1109/ICC42927.2021.9500694.

\bibitem{newman2009}
M.~E.~J.~Newman,
Random graphs with clustering,
\emph{Physical Review Letters} \textbf{103} (2009), 058701.

\bibitem{sood2024thesis}
M.~Sood,
\emph{Structural Heterogeneity and Performance in Stochastic Networks: From
Distributed Inference to Epidemics and Beyond},
PhD dissertation, Carnegie Mellon University, Pittsburgh, PA, August 2024.

\bibitem{eschenauerGligor2002}
L.~Eschenauer and V.~D.~Gligor,
A key-management scheme for distributed sensor networks,
in: \emph{Proceedings of the 9th ACM Conference on Computer and Communications
Security (CCS)}, 2002, pp.~41--47.

\end{thebibliography}

\end{document}
